\begin{table}[t]
\renewcommand{\arraystretch}{1.4}
\centering
\small
\begin{tabular}{p{0.9in} p{0.7in} p{0.8in} p{0.9in} p{0.5in} p{0.4in}}
\toprule
Mech & Loss & Mech. opt. \newline cost & Memory \newline overhead  & Noise Added & $(n,\minsep)$-fragility \\
\midrule
\blt (ours)
  & \maxloss/ \rmsloss
  & Low \newline ($\calO(n)$ )  
  & Low \newline ($\sim 4 \times \mdim$  )         
  & Low
  & Low \\
\bandmf\ifthenelse{\boolean{acl}}{}{\citep{choquette2023amplified}}
  & \rmsloss
  & High \newline ($\calO(n^2)$)
  & High \newline ($\bands \times \mdim$ )     
  & Lowest 
  & Med \\
\bandtoep\ifthenelse{\boolean{acl}}{}{\citep{mckenna2024scaling}}
  & \rmsloss
  & Low  \newline ($\calO(n)$)
  & High \newline ($\bands \times \mdim$)
  & Low 
  & Med \\
% TODO - double check on memory overhead for Full Honaker.
\treeagg\ifthenelse{\boolean{acl}}{}{\citep{kairouz21practical}}  
  & \rmsloss\!\!*
  & Lowest \newline (predefined)
  & Med \newline  ($\ceil{\log_2(\nd)} \times \mdim$) 
  & High 
  & Low  \\
\bottomrule
\end{tabular}
\caption{\small Summary of mechanisms considered, evaluated in terms of: \emph{mechanism optimization cost}, how expensive is it to compute the mechanism $\bfC$; $\calO(\cdot)$ gives the cost of a single gradient calculation. The next two columns relate to the deployment of the mechanism in a ML training system: \emph{memory overhead} is the additional state (as a multiple of the model dimension $\mdim$) that the server needs to maintain; the per-round runtime cost is also proportional to this value. The \emph{noise added} is categorized subjectively, for details see examples in ~\cref{fig:rmse}. $(n, \minsep)$-fragility reflects the degree to which the mechanisms performance degrades when total number of rounds $n$ and min-sep $\minsep$ are poorly estimated, see discussion on quantitative results in \cref{fig:rmse_k3,fig:rmse_nkb,fig:rmse_vary_n,fig:other_mechs}.  The conclusion:  \textbf{\blts perform well on all aspects.}
} \label{tab:mechanism_summary}
\end{table}


\subsection{Comparing DP-FTRL Mechanisms}
We summarize DP-FTRL mechanisms in \cref{tab:mechanism_summary} \ifthenelse{\boolean{acl}}{in \cref{app:blt-dp-ftrl}}{} and show the advantages of BLT-DP-FTRL. Our BLT mechanism can optimize either MaxLoss or RmsLoss for generating correlated noise (detailed in \cref{sec:opt_blt} following~\citep{mcmahan2024efficient}), while previous MF mechanisms in practice primarily consider RmsLoss~\citep{choquette2023amplified,mckenna2024scaling}. It is possible to extend the previous mechanisms to use MaxLoss, while it is still an open problem which loss format corresponds better with learning performance when running with the DP-FTRL algorithms. \treeagg, especially \tafull, is equivalent to considering RmsLoss though the mechanism is predefined without explicit optimization cost; and we present the lower memory overhead ($\ceil{\log_2(\nd)} \times \mdim$) for \treeagg while \tafull without an efficient algorithm yet actually needs . 

BLTs achieve better privacy-utility trade-offs than \tafull in simulation benchmark experiments (see \cref{sec:exp}), and clearly outperforms \treeagg in production cross-device FL experiments (see \cref{sec:gboard}), as lower noise is added in BLTs. While \bandmf\citep{choquette2023amplified} can add lowest noise (measured by RmsLoss in \cref{fig:rmse}), BLTs have lower mechanism optimization cost and memory overhead. Moreover, \cref{sec:exp,sec:gboard} show the learning performance of BLTs are often comparable with \bandmf under the same privacy guarantee in practical settings, though BLTs' RmsLoss is slightly worse. The memory overhead of BLTs is $d\times m$ where we empirically observe that buffer size $d=4$ achieves low losses and further increasing $d$ does not further improve in our empirical settings of total rounds $n$ and min-sep $b$. The BLT memory overhead of $d \sim 4$ is smaller than \treeagg where $\ceil{\log_2(n)} \sim 11$, and much smaller than typical $\hat b \sim X00$ for \bandmf and \bandtoep. \bandtoep~\citep{mckenna2024scaling} suggested small $\hat b$ is preferred when using amplification by sampling in the many participation settings; however, sampling is generally not possible in practical cross-device FL systems. 

As shown in \cref{fig:rmse_k3,fig:rmse_nkb,fig:rmse_vary_n,fig:other_mechs}, BLT is also more robust than banded MF mechanisms when number of total rounds $n$ and min-sep $b$ are not accurately estimated. Specifically, it is unclear how to run Banded MF mechanisms beyond the estimated $n$ after optimizing the $\bfC \in \R^{n, n}$ matrix for correlated noise. Optimizing $\bfC \in \R^{n, n}$ for a much larger $n$ and truncating it to the actual number of training rounds can achieve good privacy-utility trade-offs, but encounter non-trivial mechanism optimization cost. BandMF performance degrades fast when the actual min-sep $b$ is smaller than the estimated band $\hat b$, but the stronger DP guarantees are generally achieved when $\hat b$ is large. Hence the tricky task of estimating min-sep $b$ is more important for \bandmf. In general, BLT-DP-FTRL is competitive for achieving state-of-the-art privacy-utility trade-off (compared to \bandmf), while maintains ease to use in practice (compared to \treeagg).  